% !BIB program = biber

\documentclass{article}

% Document class and packages

\usepackage[UTF8]{ctex}  % For Chinese support

\usepackage{xeCJK}
\usepackage{amsmath,amssymb}
\usepackage{biblatex}
\addbibresource{references.bib}  % Adjust path as needed

\usepackage{listings}
\lstset{
    basicstyle=\fontspec{Consolas}\ttfamily,
    breaklines=true,
    captionpos=b,
    numbers=left,
    numberstyle=\tiny,
    frame=single,
    language=Python,  % or other languages
    literate=*{≤}{\(\leq\)}1 {≥}{\(\geq\)}1 {≠}{\(\neq\)}1 {≈}{\(\approx\)}1 {∈}{\(\in\)}1 {∉}{\(\notin\)}1 {∃}{\(\exists\)}1 {∀}{\(\forall\)}1 {⊆}{\(\subseteq\)}1 {⊂}{\(\subset\)}1 {⊇}{\(\supseteq\)}1 {⊃}{\(\supset\)}1 {∩}{\(\cap\)}1 {∪}{\(\cup\)}1 {∧}{\(\wedge\)}1 {∨}{\(\vee\)}1 {¬}{\(\neg\)}1 {⇒}{\(\implies\)}1 {⇔}{\(\iff\)}1 {∑}{\(\sum\)}1 {∏}{\(\prod\)}1 {∫}{\(\int\)}1 {∂}{\(\partial\)}1 {∇}{\(\nabla\)}1 {∞}{\(\infty\)}1 {π}{\(\pi\)}1 {μ}{\(\mu\)}1 {σ}{\(\sigma\)}1 {α}{\(\alpha\)}1 {β}{\(\beta\)}1 {γ}{\(\gamma\)}1 {δ}{\(\delta\)}1 {ε}{\(\epsilon\)}1 {θ}{\(\theta\)}1 {λ}{\(\lambda\)}1 {ρ}{\(\rho\)}1 {τ}{\(\tau\)}1 {φ}{\(\phi\)}1 {ψ}{\(\psi\)}1 {ω}{\(\omega\)}1 {Γ}{\(\Gamma\)}1 {Δ}{\(\Delta\)}1 {Θ}{\(\Theta\)}1 {Λ}{\(\Lambda\)}1 {Ξ}{\(\Xi\)}1 {Π}{\(\Pi\)}1 {Σ}{\(\Sigma\)}1 {Υ}{\(\Upsilon\)}1 {Φ}{\(\Phi\)}1 {Ψ}{\(\Psi\)}1 {Ω}{\(\Omega\)}1
}

% Other useful packages
\usepackage{graphicx}
\usepackage{float}
\usepackage{hyperref}
\hypersetup{
    colorlinks=true,
    linkcolor=blue,
    filecolor=magenta,
    urlcolor=cyan,
}
\usepackage{geometry}
\geometry{a4paper, margin=1in}

% Custom commands if needed
\newcommand{\E}{\mathbb{E}}  % Expectation
\newcommand{\cvar}{\mathrm{CVaR}}  % Conditional Value at Risk

% Title and author (can be customized in main.tex)
\title{Your Homework Title}
\author{Your Name}


\begin{document}

\maketitle

\section{引言}

在现代数学和计算机科学中，期望值（Expectation）是一个重要的概念。它描述了随机变量的平均值。在金融风险管理中，条件价值-at-风险（Conditional Value at Risk, CVaR）是一种常用的风险度量。本文将介绍CVaR的基本概念，并通过Python代码演示其计算方法。

\section{CVaR的基本概念}

CVaR是VaR的扩展，它衡量超过VaR阈值的损失的期望值。数学上，如果随机变量\(X\)表示损失，则VaR在置信水平\(\alpha\)下的定义为：
\[
VaR_\alpha = \inf \{ x \in \mathbb{R} : P(X \leq x) \geq \alpha \}
\]
而CVaR为：
\[
CVaR_\alpha = \E[ X | X > VaR_\alpha ]
\]

\section{Python实现}

以下是使用NumPy计算CVaR的简单Python代码示例：

\begin{lstlisting}[caption=CVaR计算示例,label=lst:cvar-demo]
import numpy as np

def calculate_cvar(losses, alpha=0.95):
    """
    计算CVaR
    losses: 损失数组
    alpha: 置信水平
    """
    losses = np.array(losses)
    var = np.percentile(losses, (1 - alpha) * 100)
    cvar = np.mean(losses[losses > var])
    return cvar

# 示例数据
losses = np.random.normal(100, 20, 1000)
cvar_value = calculate_cvar(losses)
print(f"CVaR at 95%: {cvar_value:.2f}")
\end{lstlisting}

\section{结论}

通过上述公式和代码，我们可以看到CVaR如何帮助量化金融风险。

\end{document}